% Ch11 WS3 -- factors affecting solubility + FRQ, plus optional enrichment.
\documentclass{apchem-worksheet}

\apcunitword{Chapter}
\apcunit{11}
\apcunittitle{Properties of Solutions}
\apcdoctitle{Worksheet 3 \textbullet\ Solubility Factors, FRQ $+$ Optional}
\apcsubtitle{Zumdahl \S11.3 examinable \textbullet\ \S11.4--11.7 optional}

\begin{document}
\wsheader[Problems 1--4 are exam-relevant. Problem 5 is marked OPTIONAL ---
colligative properties are not on the AP exam.]

\problem[]{Predict which is more soluble in water and give the structural
reason:}
\begin{parts}
  \item \ce{CH3CH2OH} or \ce{CH3(CH2)6CH2OH}
        \answerlines[2]{Ethanol --- both have one O--H, but the longer chain
        adds nonpolar bulk that water cannot solvate, so the polar fraction
        of the molecule shrinks.}
  \item \ce{NH3} or \ce{PH3}
        \answerlines[2]{\ce{NH3} --- it hydrogen bonds with water; \ce{PH3}
        cannot, since phosphorus is not electronegative enough.}
  \item \ce{KNO3} or \ce{I2}
        \answerlines[2]{\ce{KNO3} --- its ions are stabilized by strong
        ion--dipole attractions; \ce{I2} is nonpolar.}
\end{parts}

\problem[]{Vitamins A and C.}
\begin{parts}
  \item Vitamin A is largely carbon and hydrogen; vitamin C carries several
        O--H and C--O bonds. Classify each as hydrophobic or hydrophilic.
        \answerblank[6cm]{A hydrophobic; C hydrophilic}
  \item Explain why excess vitamin A can cause hypervitaminosis while excess
        vitamin C generally does not. \answerlines[3]{Vitamin A dissolves in
        fatty tissue and accumulates there, so an excess builds up. Vitamin
        C stays in the aqueous phase and is excreted, so it does not
        accumulate --- but it must be consumed regularly.}
\end{parts}

\problem[]{Pressure and temperature effects. Complete the table:}

\renewcommand{\arraystretch}{1.5}
\noindent\begin{tabular}{@{}p{3.2cm}p{6.0cm}p{5.0cm}@{}}
\toprule
 & \textbf{Solids in water} & \textbf{Gases in water}\\
\midrule
Increasing pressure & \answerblank[4.1cm]{essentially no effect} &
  \answerblank[3.9cm]{increases solubility}\\
Increasing temperature & \answerblank[5.8cm]{usually increases, not always} &
  \answerblank[4.0cm]{decreases solubility}\\
\bottomrule
\end{tabular}

\problem[]{Explain each observation:}
\begin{parts}
  \item An opened bottle of soda goes flat.
        \answerlines[2]{Sealed, the \ce{CO2} is under high partial pressure
        and much dissolves. Opened, it equilibrates with the atmosphere
        where the \ce{CO2} partial pressure is minuscule, so nearly all the
        dissolved gas escapes.}
  \item A warm soda fizzes more violently than a cold one when opened.
        \answerlines[2]{Gas solubility falls as temperature rises, so the
        warm liquid can hold less \ce{CO2} at equilibrium and expels it
        faster.}
  \item A power plant discharging warm water can suffocate fish.
        \answerlines[2]{Warmer water holds less dissolved oxygen, and
        aquatic organisms depend on that dissolved \ce{O2}.}
  \item A student claims heating always dissolves more solute. Correct them
        precisely. \answerlines[2]{Heating always makes dissolving
        \emph{faster}, but the maximum amount that dissolves may increase or
        decrease --- sodium sulfate becomes less soluble as temperature
        rises.}
\end{parts}

\problem[]{\textbf{FRQ (10 points).} A student investigates the dissolving
of two ionic salts in water. Salt A causes the solution temperature to fall
by \SI{8.2}{\celsius}; salt B causes it to rise by \SI{11.5}{\celsius}.}
\begin{parts}
  \item State the sign of $\Delta H_{\text{soln}}$ for each salt and name
        the process. \answerlines[2]{Salt A: positive, endothermic --- it
        absorbed energy from the solution. Salt B: negative, exothermic.}
  \item Using the three-step model, explain how salt A can dissolve at all
        despite absorbing energy. \answerlines[3]{Steps 1 and 2 (breaking
        the lattice and expanding the solvent) together exceed the energy
        released by Step 3 (hydration), so $\Delta H_{\text{soln}} > 0$.
        Dissolving still occurs because the increase in disorder on mixing
        can drive the process even when it is energetically uphill.}
  \item For salt B, state which energy term must be unusually large, and
        identify it by name. \answerlines[2]{Step 3 --- the hydration
        energy. It must exceed the combined cost of breaking the lattice and
        expanding the solvent.}
  \item Which salt would be suitable for an instant cold pack, and which for
        a hand warmer? \answerblank[6cm]{A cold pack; B hand warmer}
  \item The student dissolves \SI{0.150}{\mole} of salt A in
        \SI{200.0}{\gram} of water. Calculate $\Delta H_{\text{soln}}$ in
        \si{\kilo\joule\per\mole}, assuming the solution has the specific
        heat of water (\SI{4.184}{\joule\per\gram\per\celsius}).
        \workspace[2.6cm]
        \keyonly{\begin{apcanswer}$q_{\text{soln}} = (200.0)(4.184)(-8.2) =
        \SI{-6862}{\joule}$, so the dissolving absorbed
        $+6.86$ kJ. $\Delta H_{\text{soln}} = 6.86/0.150 =
        \SI{+45.7}{\kilo\joule\per\mole}$\end{apcanswer}}
\end{parts}

\begin{apcnote}[Rubric --- score yourself, 10 points]
\textbf{(a) 2 pts} 1 per salt, sign AND term \textbullet\ \textbf{(b) 3 pts}
1 identifies steps 1+2 exceeding step 3, 1 states $\Delta H > 0$,
1 invokes disorder/entropy as the driver \textbullet\ \textbf{(c) 2 pts}
1 names step 3, 1 calls it hydration energy \textbullet\ \textbf{(d) 1 pt}
both assignments \textbullet\ \textbf{(e) 2 pts} 1 flips the sign from
solution to system, 1 correct value per mole.
\end{apcnote}

\problem[]{\textbf{OPTIONAL --- enrichment only.} Colligative properties are
not assessed on the AP exam.}
\begin{parts}
  \item Define a colligative property.
        \answerlines[2]{One that depends only on the number of dissolved
        particles, not on their chemical identity.}
  \item Rank by freezing-point depression at equal molality: glucose,
        \ce{NaCl}, \ce{CaCl2}. Explain.
        \answerlines[2]{\ce{CaCl2} $>$ \ce{NaCl} $>$ glucose --- they
        release 3, 2, and 1 particles per formula unit respectively.}
  \item Explain why a red blood cell bursts in pure water.
        \answerlines[2]{Water moves by osmosis from the dilute exterior into
        the more concentrated cell interior, swelling it until the membrane
        fails.}
\end{parts}

\begin{spiralstrip}{Chapter 11 \textbullet\ blocks 1--2}
  \item Molality formula: \answerblank[4.5cm]{mol solute / kg solvent}
  \item Sign of $\Delta H$ for expanding the solvent:
        \answerblank[2.2cm]{positive}
  \item Cold-pack salt: \answerblank[2.6cm]{\ce{NH4NO3}}
  \item For an ionic solute, step 3 is called:
        \answerblank[3.4cm]{hydration energy}
  \item Which is temperature-independent, molarity or molality?
        \answerblank[2.4cm]{molality}
\end{spiralstrip}

\end{document}
